\scanpage{42}
% L1–4
\uline{Proposition} Supposons le foncteur $N\to\operatorname{Ext}^{i}(N,M)$ ($N$ de longueur finie) \uline{exact \unsure{à ga|à g.}} Alors on a
\[
 \operatorname{Ext}^{i}(N\deleted{\unsure{,}},M)\simeq\operatorname{Hom}(N,H^{i}(M))
\]
\margin{[\deleted{a-t-on} \deleted{\unsure{alors $i$|alors $i=$|donc} \illegible{1}} $\operatorname{Ext}^{i}(k,M)$ ou $\operatorname{Ext}^{i-1}(k,M)=0$ ??]}

% L5–11
Supposons le foncteur \uline{exact}, alors ou bien il est identiquement nul, [i.e. \deleted{$H^{i}(M)=0$, \unsure{ou c'est|ou encore}} $\operatorname{Ext}^{i}(k,M)=0$, cela implique prob.$^{\mathrm{t}}$ $i>n$ \uuline{ou} $i<\deleted{\unsure{d|co}}\dim.\,\mathrm{coh.}\ M$], ou bien $i=n$, $\widehat{H^n(M)}\simeq A^{m}$. Alors $M$ est \uline{fidèle} $\operatorname{Ext}^{n}(k,M)\simeq k^{m}$. De plus, cela implique $\operatorname{Ext}^{i}(k,M)=0$ pr $i>n$.
\margin{[et \unsure{cela|ceci} implique-t-il que $A$ est \uline{Coh. Mac.}, et $M\simeq\Omega_A^m$] \deleted{\unsure{fidélité|fidèle}}}

\par\noindent\shortstack{\rule{.4\linewidth}{.6pt}\\[-1ex]\rule{.4\linewidth}{.3pt}}

% L12–16
% slanted
\margin{$M$ un \unsure{mod.|module|néc.} de type fini}
Supposons le foncteur $N\mapsto\operatorname{Ext}^{\deleted{\unsure{i|i+1}}i-1}(N,M)$ \uline{exact} \uuline{à droite}. Je dis qu'on a \deleted{ou bien} $\operatorname{Ext}^{\deleted{\unsure{i+1|i}}\unsure{i-1|i}}(k,M)=0$, ou bien $\operatorname{Ext}^{\deleted{\unsure{i+1|i-1}}i}(k,M)=0\Longleftrightarrow H^{i}(M)=0$.

En effet, on aura $\operatorname{Ext}^{i}(N,M)\simeq\operatorname{Hom}(N,H^{i}(M))$.

% L17
\diagram{\begin{tikzcd}[ampersand replacement=\&]
0 \arrow[r] \& M \arrow[r] \& J \arrow[r] \& M' \arrow[r] \& 0
\end{tikzcd}}

% archive sheet number: 41
